\scanpage{213}
\uline{Groupe fondamental motivique discret}

Soit $\xi\in\mathbf T(\mathbf Z)$, on pose
\[
\begin{aligned}
\pi_1(X,\xi)&=\operatorname{Aut}(\xi).\\
&=G_\xi(\mathbf Z)
\end{aligned}
\]

\uline{N.B.} En vertu de Borel\inserted{-\unsure{Housh|Harish}}, il n'y a qu'un nb fini de $\xi\in T_\alpha(\mathbf Z)$ possible, pour $\alpha$ donné.
D'où un nb fini de \uline{groupes fondamentaux discrets}.
Sont-ils tous \unsure{associés} à des \deleted{\unsure{champs}} \unsure{places} ??

$\xi\in T_\alpha(k)$, $k$ un anneau

$M$ motif $\in\operatorname{Ob}\unsure{\mathfrak M_\alpha|\mathcal M_\alpha}$

$M_\xi$ fibre en $\xi$,\qquad $M_\xi=\xi(M)$
\[
\pi_1(X,\xi)_{\unsure{a|\alpha}}=\operatorname{Aut}(\xi),\qquad
G(X,\xi)=\underline{\operatorname{Aut}}(\xi)
\]
\deleted{$\unsure{G_\xi}$}
\[
\pi_1(X,\xi)_{\unsure{a|\alpha}}\simeq G(X,\xi)(k)
\]
Comparaison avec $\pi_1(X^{\unsure{\mathrm{tr}|\mathrm{tv}}},\unsure{\bar x|\xi})$ quand $\xi$ provient d'un \uline{pt} \unsure{en $a$} ....

\uline{Ex}. $X$ connexe sur $\mathbf C$, $\xi$ défini par un pt de $X$ à valeurs dans $\mathbf C$.
Le foncteur $\xi_{\mathbf C}\deleted{\unsure{\otimes}}$ a une structure \uline{filtrée}.
$\xi$ considéré comme foncteur à valeurs dans les modules sous $\pi_1(X,\xi)$ muni d'une \uline{structure de Hodge} est \unsure{pl.} fidèle ....

\margin{212}
